%======================================================================
\section{Versión 7 y reconstrucción v8-r}\label{sec:v8r}

\paragraph{Estado.} El código fuente de la versión 8 se perdió. La implementación
actual (\texttt{sotsort.hpp}, ``v8-r'') se reescribió a partir de este documento, del
libro y de una versión anterior del código. Las demostraciones de las secciones
anteriores son sobre el pseudocódigo y no cambian; lo que cambia es que la
correspondencia con el código vuelve a apoyarse solo en pruebas diferenciales contra
\texttt{std::stable\_sort} (sección~\ref{sec:v8r-pruebas}). Esta sección agrega lo que
el documento no cubría: dos cambios de la versión 7 y cuatro adiciones de v8-r. Las
cuatro adiciones solo eligen entre procedimientos ya demostrados o reordenan
recorridos; ninguna introduce un procedimiento de ordenamiento nuevo.

\subsection{Versión 7: pivote con posiciones pseudoaleatorias}
\Pivote toma sus muestras en posiciones $l+js+\Aleatorio{}\bmod s$ en lugar de $l+js$,
con un generador determinista sembrado con $\sigma\oplus n$. El Lema~\ref{lem:pivote}
solo usa que $p$ es copia de un elemento de \seg{A}{l}{r}, y cada posición muestreada
cumple $l\le l+js+(\Aleatorio{}\bmod s)<r$ porque $(j+1)s\le r-l$. El lema y todo lo que
depende de él valen sin cambios. El Teorema~\ref{thm:det} tampoco cambia: la secuencia de
\Aleatorio es función de $n$ y de la semilla.

\subsection{Versión 7: partición de bloque vectorial}

\begin{algorithm}[H]
\caption{Partición de bloque vectorial, con $w$ carriles}\label{alg:partvec}
\KwIn{$s<e$ con $L=e-s\le\Bsz$; pivote $p$; $\triangleleft\in\{\lt,\les\}$; $T$ entero y $\lt$ es $<$}
\Fn{\PartVec{$s, e, p, \triangleleft$}}{
  $m\gets 0$;\quad $i\gets 0$\;
  \While{$i+w\le L$}{
    $X\gets A[s+i\,..\,s+i+w)$\tcp*{una carga}
    $k\gets\{\,j<w : X[j]\triangleleft p\,\}$\;
    $A[s+m\,..\,s+m+w)\gets\mathrm{comprimir}(X,k)$\;
    $\buf[i-m\,..\,i-m+w)\gets\mathrm{comprimir}(X,\overline{k})$\;
    $m\gets m+|k|$;\quad $i\gets i+w$\;
  }
  cola escalar: el lazo de \PartBloque desde $i$\;
  $A[s+m\,..\,e)\gets\buf[0\,..\,L-m)$;\quad \KwRet $s+m$\;
}
\end{algorithm}

\noindent $\mathrm{comprimir}(X,k)$ es el vector con los carriles de $X$ marcados en $k$,
en su orden, seguidos de $w-|k|$ carriles de contenido arbitrario.

\begin{lema}[Partición vectorial]\label{lem:partvec}
Con $C=\seg{A_0}{s}{e}$ el contenido inicial y $P(x)\equiv x\triangleleft p$, al terminar
$\seg{A}{s}{e}=\filt(P,C)\cat\filt(\neg P,C)$ y se devuelve $s+|\filt(P,C)|$. Ninguna
escritura cae fuera de \seg{A}{s}{e} ni de $\buf[0..L)$.
\end{lema}
\begin{proof}
Invariante del lazo: $\seg{A}{s}{s+m}=\filt(P,C[0..i))$,
$\buf[0..i-m)=\filt(\neg P,C[0..i))$ y $\seg{A}{s+i}{e}=C[i..L)$.
\emph{Preservación.} $X=C[i..i+w)$ por la tercera cláusula. El store sobre $A$ escribe
$[s+m,s+m+w)$; como $m\le i$, ese rango está contenido en $[s,s+i+w)$, cuyo contenido
ya fue leído ($C[0..i)$ en vueltas anteriores, $C[i..i+w)$ en $X$), así que no destruye
nada de $C[i+w..L)$, y $s+m+w\le s+i+w\le e$. Sus primeros $|k|$ carriles son
$\filt(P,X)$ y extienden el prefijo. El store sobre $\buf$ escribe
$[i-m,i-m+w)$ con $i-m+w\le i+w\le L\le\Bsz$; sus primeros $w-|k|$ carriles son
$\filt(\neg P,X)$. \emph{Terminación.} La cola escalar es el lazo de
\PartBloque con el mismo invariante (Lema~\ref{lem:partbloque}), y la copia final es
la suya.
\end{proof}

Por el Supuesto~\ref{sup:enteros}, evaluar $\triangleleft$ con una comparación vectorial
de enteros da el mismo resultado que $\lt$ carril por carril. El
Lema~\ref{lem:partest} usa a \PartBloque solo a través de su poscondición, que
\PartVec comparte.

\subsection{v8-r: sonda de rachas en SymMerge}\label{sec:sonda}
En una hoja de \SymMerge (un lado entra en el buffer) hay tres mezclas posibles: copiar
todo al buffer y hacer un \Parity de largos distintos, si los dos runs entran juntos;
copiar el lado corto y usar \MIzq/\MDer con galope; o \MIzq/\MDer con el lazo
clásico. La primera no tiene saltos y conviene cuando los runs se intercalan elemento
por elemento; el galope conviene con rachas largas. En el camino de runs naturales,
para tipos que se copian byte a byte y $a+b\ge256$, se decide con una sonda.

Sean $X=\seg{A}{l}{m}$ e $Y=\seg{A}{m}{r}$ ordenados, $a=|X|$, $b=|Y|$. La sonda toma
$x_0=l$, $x_i=l+\lfloor ia/9\rfloor$ para $1\le i\le 8$, $x_9=m$, y
$y_0=m$, $y_i=\min\{\,j\in[y_{i-1},r] : j=r \vee A[j]\nlt A[x_i]\,\}$ para $1\le i\le8$,
$y_9=r$. Define
\[\mathit{sw}=\sum_{i=1}^{9}\bigl(2\min(x_i-x_{i-1},\,y_i-y_{i-1})+1\bigr).\]

\begin{lema}[Sonda de rachas]\label{lem:sonda}
\begin{enumerate}[label=(\alph*),nosep]
\item $\Mrg(X,Y)$ tiene a lo sumo $\mathit{sw}$ cambios de lado (posiciones consecutivas
  de la salida que vienen una de $X$ y otra de $Y$).
\item La sonda hace a lo sumo $8\lceil\lg(b+1)\rceil$ comparaciones y no escribe.
\item Elegir una u otra mezcla según $\mathit{sw}$ no altera la poscondición del
  Lema~\ref{lem:sym}.
\end{enumerate}
\end{lema}
\begin{proof}
(a) Para $1\le i\le8$: todo $y$ de $Y$ antes de $y_i$ cumple $y\lt A[x_i]$ (los $y_i$ son
no decrecientes e $Y$ está ordenado), así que precede a $A[x_i]$ en $\Mrg(X,Y)$; todo $y$
desde $y_i$ cumple $y\nlt A[x_i]$ y, como en empates va primero $X$
(Def.~\ref{def:merge}), queda después de $A[x_i]$. Luego
$\Mrg(X,Y)=\cat_{i=1}^{9}\Mrg(\seg{A}{x_{i-1}}{x_i},\seg{A}{y_{i-1}}{y_i})$. Una mezcla
de $g_x$ y $g_y$ elementos tiene a lo sumo $2\min(g_x,g_y)$ cambios de lado, y entre dos
tramos consecutivos hay a lo sumo uno más.
(b) Son 8 búsquedas binarias sobre sufijos de $Y$.
(c) Las tres mezclas cumplen la misma poscondición: Lemas~\ref{lem:parity},
\ref{lem:mizq} y~\ref{lem:mizqg}.
\end{proof}

La regla es: galope si $8\,\mathit{sw}\le a+b$ (por (a), las rachas miden 8 o más en
promedio); si no, \Parity cuando $a+b\le\Bsz$, y el lazo clásico con saltos cuando no
entran juntos. Con menos de 256 elementos se usa \Parity sin sondear. El umbral y la
cantidad de muestras se eligieron midiendo en una sola máquina. Para los tipos que no se
copian byte a byte no hay sonda: se usa el galope, como en la versión 8. Por (b) y
porque \Natural hace a lo sumo $\lceil n/24\rceil-1$ mezclas, la sonda agrega
$O(n\log n)$ comparaciones en el peor caso y no cambia el Teorema~\ref{thm:cost}; con
mezclas balanceadas de tamaño $s\ge256$ hay $n/s$ por nivel y el total es $O(n)$.

\subsection{v8-r: sonda de orden global y regla de ruta}\label{sec:ruta}
\PocosGiros manda a \Natural toda entrada con menos de $n/8$ giros. Eso es un error
cuando los runs tienen claves entremezcladas (por ejemplo, tramos ordenados de 63 claves
al azar): \Natural hace $\lg(\text{runs})$ niveles de mezclas que alternan elemento por
elemento. En v8-r, si hay pocos giros se toman 256 muestras, una por tramo de largo
$\lfloor n/256\rfloor$ en una posición pseudoaleatoria del tramo, y se cuentan los
cambios de dirección entre muestras vecinas. Con menos de 64 cambios la entrada tiene
orden global y va a \Natural. Si no, va a \Natural solo cuando
$k\lceil\lg R\rceil\le\lceil\lg(n/\leaf)\rceil$, con $R=\lfloor g/2\rfloor+1$ la
estimación de runs a partir de los $g$ giros y $k$ el costo relativo de un nivel de
mezcla ($k=3$ o $2$ con buffer chico, según haya o no partición vectorial; $2$ o $5/4$
con buffer de al menos $n/2$). La sonda y la regla se aplican solo a tipos que se copian
byte a byte: con \emph{strings} mandar esas entradas al quicksort midió peor.

\begin{observacion}\label{obs:ruta}
La regla solo elige entre \Natural y \Quick. Los Teoremas~\ref{thm:natural}
y~\ref{thm:quick} valen para toda entrada que cumple sus precondiciones, sin importar
cómo se eligió el camino, así que el Teorema~\ref{thm:main} no cambia. La sonda hace
255 comparaciones y no escribe. Las constantes $k$ y el umbral de 64 salen de mediciones
en una sola máquina; una mala elección cuesta tiempo, nunca corrección. Sin el
desplazamiento pseudoaleatorio, runs de largo fijo producen muestras en posiciones
relativas que avanzan en progresión aritmética, y la entrada parece ordenada
(lo vi con runs de 1000 y $n=10^6$).
\end{observacion}

\subsection{v8-r: zonas desordenadas en MergeNatural}\label{sec:zonas}
Una entrada con orden global puede tener una zona al azar (el 10\% final, o un bloque
en el medio). Ahí \Natural encuentra runs de uno o dos elementos, los extiende a 24 con
inserción y los mezcla con mezclas mal predichas. En v8-r, cuando \Natural encuentra un
run de menos de 24 elementos mira hacia adelante: si ese run y los siguientes suman 8
o más runs seguidos de menos de 8 elementos, extiende la zona $[i,e)$ hasta el primer
run de 16 o más (o el final), la ordena con \Quick y la apila como un solo run. Si no,
hace lo de antes. Igual que la regla de ruta, se aplica solo a tipos que se copian byte
a byte.

\begin{lema}[Zonas]\label{lem:zonas}
Con esa modificación, el Teorema~\ref{thm:natural} sigue valiendo, y la pila de runs y
la cota de mezclas por elemento del Lema~\ref{lem:natcost} no cambian. El recorrido
hacia adelante agrega a lo sumo $3n$ comparaciones en total.
\end{lema}
\begin{proof}
El invariante de \Natural pide que cada entrada de la pila sea un intervalo ordenado y
coherente con los elementos de $\Omega$ restringido a él, y que los intervalos sean
adyacentes. La zona $[i,e)$ es un intervalo todavía no apilado; su contenido es
coherente porque solo se invirtieron runs estrictamente descendentes
(Lema~\ref{lem:yaord}). Por el Teorema~\ref{thm:quick}, \Quick deja $[i,e)$ ordenado y
coherente, que es lo que el invariante pide de un run. El buffer está libre entre dos
mezclas, así que \Quick puede usarlo. El Lema~\ref{lem:natcost} solo usa los largos de
los runs apilados, no cómo se obtuvieron, y la zona mide al menos 8 elementos. Costo
del recorrido: cada run mirado se recorre hasta 16 elementos como máximo, se miran a lo
sumo 8 runs antes de decidir que no hay zona (menos de $7\cdot 8+16=72$ comparaciones), y
después de decidirlo \Natural avanza al menos 24 posiciones; cuando hay zona, el
recorrido es lineal en la zona.
\end{proof}

El tiempo de la zona es el de \Quick sobre $e-i$ elementos; el Teorema~\ref{thm:cost}
no cambia porque \Quick está dentro de la misma cota.

\subsection{v8-r: recorridos vectoriales}
Para enteros con el orden natural, \Ordenado y \PocosGiros comparan $w$ pares adyacentes
por instrucción. Calculan exactamente los mismos predicados
($A[i]\lt A[i-1]$ para cada $i$) en el mismo orden de posiciones, así que los
Lemas~\ref{lem:yaord} y la Observación~\ref{obs:heur} no cambian. Las lecturas van de
$A[i-1]$ a $A[i+w-1]$ con $i+w\le r$: no salen del rango.

\subsection{Verificación de v8-r}\label{sec:v8r-pruebas}
Pruebas diferenciales contra \texttt{std::stable\_sort}, con el índice original como
etiqueta: enteros de 32 y 64 bits con y sin signo, pares, registros de 16, 40 y 128
bytes y \emph{strings}; los tres modos de memoria; de 1 a 4 hilos; orden por clave;
15 familias de entradas, incluidas zonas desordenadas y runs de claves al azar. Sobre el
archivo final: sin fallas en 19\,506 casos con AVX-512 y en 7\,410 con AVX2, sin SIMD y
con clang. AddressSanitizer, UBSan y ThreadSanitizer no informan nada. Se inyectaron
fallas de reserva en cada punto. Un contador en una versión anterior del mismo día
confirmó que las pruebas ejecutan las zonas (207 veces en una corrida), la regla de ruta
(480) y la sonda de orden (1245). Limitación conocida: si el
comparador lanza una excepción se pueden perder elementos; no hay errores de memoria.

%======================================================================
\section{Planificación fuera de orden (experimental)}\label{sec:ooo}

Esta sección estudia ejecutar los rangos pendientes de \Quick como tareas, en cualquier
orden y en varios hilos, sin las rondas de mezcla de \Paralelo. El prototipo
(\texttt{sotsort\_ooo.hpp}) está aislado de la biblioteca. Los resultados (a) a (e) del
Teorema~\ref{thm:ooo} están demostrados a mano al nivel de este documento; la
implementación del planificador solo está verificada por pruebas.

\subsection{Modelo}
Una \emph{tarea} es un rango pendiente $(l,r,\beta,\mathit{lb})$ de \Quick. Ejecutarla
es correr el lazo interior del Algoritmo~\ref{alg:quick} sobre ese rango, con una
diferencia: cuando el algoritmo secuencial apilaría la mitad grande de una partición,
la tarea puede \emph{cederla}: la pone en una cola compartida y sigue con la mitad
chica. La cola admite a lo sumo $64P$ tareas; con la cola llena la mitad grande se
apila localmente, como en el algoritmo secuencial. $P$ trabajadores toman tareas de la
cola hasta que está vacía y ninguno trabaja. Ninguna tarea espera a otra.

La \emph{huella} de una tarea es el conjunto de celdas que lee o escribe. Dos tareas
\emph{conmutan} si ejecutarlas en cualquier orden, o intercaladas, deja el mismo estado.

\begin{lema}[Conmutación]\label{lem:conm}
Si las huellas de $T_1$ y $T_2$ son disjuntas, $T_1$ y $T_2$ conmutan, y bajo el
Supuesto~\ref{sup:hilos} pueden ejecutarse a la vez sin carreras.
\end{lema}
\begin{proof}
Es el caso más simple de las condiciones de Bernstein: ninguna escribe lo que la otra
lee o escribe, así que cada una observa el mismo estado inicial en cualquier
intercalado y sus efectos se componen.
\end{proof}

\begin{lema}[Buffers disjuntos]\label{lem:bufdisj}
En modo completo, si mientras \Quick procesa $[l,r)$ se usa $\buf=\mathit{aux}+l$ y
$\Bsz=r-l$, toda escritura en el buffer cae en $\mathit{aux}[l..r)$. En los otros modos,
con un buffer por trabajador, toda escritura en el buffer cae en el buffer del
trabajador que ejecuta la tarea.
\end{lema}
\begin{proof}
Segundo caso: inmediato. Primer caso, por lo que se ejecuta con ese rango.
\Part: como $r-l\le\Bsz$ hay un solo bloque; \PartBloque escribe $\buf[i-m]$ con
$i<r-l$ y \PartVec escribe $\buf[i-m..i-m+w)$ con $i+w\le r-l$
(Lema~\ref{lem:partvec}); no hay fusiones, así que \Rotar no se ejecuta. \Hoja: escribe
$\buf[0..c)$ y $\buf[0..t)$ con $c\le t\le r-l$. \Pivote: a lo sumo $\Bsz$ muestras.
\Respaldo: \SymMerge sobre un subrango $[l',r')$ escribe desde $\buf[0]$ a lo sumo
$r'-l'\le r-l$ celdas.
\end{proof}

\begin{teorema}[Planificación fuera de orden]\label{thm:ooo}
Sea $A_0$ coherente con $n\ge2$ y supongamos los
Supuestos~\ref{sup:orden}--\ref{sup:hilos}. Para cualquier orden de ejecución de las
tareas y cualquier $P\ge1$:
\begin{enumerate}[label=(\alph*),nosep]
\item en todo momento, los rangos de la cola, los de las pilas locales y los que están
  en curso son disjuntos dos a dos y cada uno es pendiente en el sentido de la
  Def.~\ref{def:invq}; el resto de $[0,n)$ está terminado;
\item dos tareas en curso tienen huellas disjuntas;
\item las dependencias entre tareas forman un árbol;
\item la ejecución termina, sin bloqueo mutuo ni inanición;
\item al terminar, $A=\Omega$: la salida no depende del orden de ejecución ni de $P$, y
  es la de \SotSort.
\end{enumerate}
\end{teorema}
\begin{proof}
(a) El invariante de \Quick (Def.~\ref{def:invq}) habla de un conjunto de intervalos
pendientes; la demostración del Teorema~\ref{thm:quick} nunca usa que estén en una
pila ni el orden en que se procesan. Cada paso de una tarea es un paso del lazo
secuencial sobre su propio rango y preserva el invariante por la misma prueba. Ceder
una mitad en lugar de apilarla no cambia el conjunto.
(b) Por (a) los rangos son disjuntos; por el Lema~\ref{lem:bufdisj} los buffers
también. El estado del planificador se lee y escribe con un cerrojo tomado, y la cota
inferior viaja por copia (Supuesto~\ref{sup:ident}).
(c) Una tarea depende solo de la partición que la creó, y su rango está contenido
estrictamente en el de esa partición.
(d) Ninguna tarea espera a otra, así que no hay ciclo de espera. Cada tarea termina
por el Lema~\ref{lem:prof}. Las tareas son finitas: los rangos cedidos son disjuntos o
anidados y cada uno mide al menos el grano. La cola se atiende en orden de llegada.
Un trabajador termina solo cuando la cola está vacía y ningún otro trabaja, y en ese
estado no puede aparecer otra tarea.
(e) Al terminar no quedan pendientes, así que por (a) todo $[0,n)$ está terminado:
$A=\Omega$. $\Omega$ es única (Def.~\ref{def:omega}).
\end{proof}

\begin{observacion}
(e) es unicidad de la permutación estable, no una propiedad del planificador: cualquier
ordenamiento estable correcto da la misma salida. Lo que sí depende del diseño es que
las comparaciones realizadas no dependan del orden de ejecución: la semilla de
\Aleatorio de cada tarea cedida se deriva de $(\sigma,l,r)$, y la decisión de ceder
depende solo de los tamaños. Si la cola se llena, la decisión pasa a depender del
orden de ejecución; en las mediciones no se llenó nunca.
\end{observacion}

\subsection{Partición paralela de la raíz}
La primera partición es secuencial y fija el camino crítico. En modo completo se
reemplaza por tres fases. Sean $d_0=0<d_1<\dots<d_P=n$ y $P(x)\equiv x\lt p$.
\textbf{Fase 1} (en paralelo, tramo $i$): los que cumplen $P$ quedan en
$A[d_i..d_i+m_i)$ y el resto en $\mathit{aux}[d_i..d_i+f_i)$, en su orden (es
\PartBloque sin la copia final). \textbf{Fase 2} (secuencial, $i=1,\dots,P-1$):
$A[d_i..d_i+m_i)$ se mueve a $A[M_i..)$ con $M_i=\sum_{j<i}m_j$. \textbf{Fase 3} (en
paralelo): $\mathit{aux}[d_i..d_i+f_i)$ se copia a $A[M+F_i..)$, con $M=\sum_j m_j$ y
$F_i=\sum_{j<i}f_j$.

\begin{lema}[Raíz paralela]\label{lem:raiz}
Al terminar, $A=\filt(P,A_0)\cat\filt(\neg P,A_0)$, y las fases 1 y 3 no tienen
carreras.
\end{lema}
\begin{proof}
Fase 1: el tramo $i$ usa $A[d_i..d_{i+1})$ y $\mathit{aux}[d_i..d_{i+1})$, disjuntos
entre tramos. Fase 2: $M_i\le d_i$, así que el destino empieza a la izquierda del origen
y un movimiento hacia adelante es seguro; el destino termina en
$M_i+m_i\le d_i+m_i\le d_{i+1}$, antes del origen del tramo siguiente. Queda
$A[0..M)=\filt(P,A_0)$ porque cada tramo conserva su orden y los tramos se concatenan en
orden. Fase 3: los destinos $[M+F_i,M+F_i+f_i)$ son disjuntos, están en $A[M..n)$, y los
orígenes están en $\mathit{aux}$; $A[M..n)$ ya no tiene nada pendiente.
\end{proof}

\subsection{Trabajo y camino crítico}
El trabajo $W(n)$ es el de \Quick más $O(1)$ por tarea. Con la partición secuencial, el
camino crítico es $S(n)=\Theta(n)$ esperado (la partición de un rango de largo $\lambda$
cuesta $\Theta(\lambda)$ y, siguiendo la mitad grande, las $\lambda$ decrecen
geométricamente), así que el paralelismo $W/S$ es $\Theta(\log n)$. Con la raíz paralela
el primer término baja a $n/P+(M-m_0)+\max_i f_i$. Dentro de \Part por bloques, con
rotaciones secuenciales, la fusión de la raíz mueve $\Theta(\lambda)$: el paralelismo
interno es $O(\lg(\lambda/\Bsz))$, y por eso no se construyó una variante con tareas por
bloque.

\subsection{Qué no está demostrado}
\begin{itemize}[nosep]
\item Que la implementación del planificador (cola, cerrojo, variable de condición)
  cumple el modelo. Evidencia: 466 casos contra \texttt{std::stable\_sort} con 2 a 5
  hilos y grano de 64 a 32\,768, y 60 corridas con un comparador que lanza, todo sin
  avisos de ThreadSanitizer ni de AddressSanitizer.
\item Cualquier afirmación de escalabilidad: solo se midió con 2 núcleos.
\item El modelo de memoria de C++ en detalle; el Supuesto~\ref{sup:hilos} lo resume.
\end{itemize}
